\chapter{Experiments}\label{chapter:experiments}

\section{Experimental Setup}

This section describes the experimental setup used to evaluate the ALERT framework. Section \ref{sec:dataset} summarizes the SoccerNet Player-Tracking dataset. Section \ref{sec:evaluation_metrics} introduces the evaluation metrics. Section \ref{sec:baselines} describes the baselines against which ALERT is compared. Section \ref{sec:training_and_evaluation} lays out the training and evaluation details, and Section \ref{sec:implementation_details} covers implementation details.

\subsection{Implementation Details}\label{sec:implementation_details}
Table \ref{tab:system_configs} show the system configuration for all experiments, including hardware and software specifications. 


\begin{table}[h]
\centering
\begin{tabular}{ll}
\toprule
\textbf{Parameter} & \textbf{Value} \\
\midrule
Processor & Ryzen 7 3700X \\
Memory & 16GB RAM \\
GPU & RTX 2070 8 GB VRAM \\
Operating System & Windows 11 \\
Python version & 3.10.18 \\

\bottomrule

\end{tabular}
\caption{System configuration details}
\label{tab:system_configs}
\end{table}

\subsection{Dataset}\label{sec:dataset}
All experiments are performed on the newly constructed SoccerNet Player-Tracking dataset (SNPT), described in detail in Section \ref{dataset_construction}. Each sequence contains 30 seconds of broadcast  football footage recorder at 25 frames per second with a resolution of $1920 \times 1080$. The dataset only contains annotations of players and goalkeepers, while referee, pitch and the ball are removed during the dataset conversion. The dataset is divided into a train, validation, and test split containing 57, 58, and 49 sequences respectively. 

\

The train split is used to train the learning-based mergers (XGBoost and transformer), while the validation split is used for hyperparameter tuning. The test split is held out and is used exclusively to report the final results, preventing data leakage and overfitting to the test set.




\subsection{Evaluation Metrics}\label{sec:evaluation_metrics}

\subsubsection{Tracklet Splitter}
Evaluating a tracklet splitter is challenging due to the way the HOTA metric works. While successfully splitting a true identity switch into pure fragments theoretically improves association accuracy, the metric simultaneously applies a harsh penalty for creating additional fragments. Because this fragmentation penalty can outweigh the reward for correct splits, a model that successfully catches many identity switches but also produces false splits will obtain a low score. As a result, the HOTA score on  the splitter output alone cannot distinguish between a precise splitter and one that avoids splitting completely. To better isolate the splitter performance, we evaluate its intrinsic metrics.


The intrinsic metrics measure the splitting quality against the ground truth identity labels without a downstream tracklet merger. The ground truth identity switches are defined as all positions inside a  tracklet where the ground truth identity changes between two consecutive detections. The baseline tracklets contain \textbf{1048} identity switches across the 49 test sequences.

\

To evaluate the tracklet splitter's decisions, we report three specific metrics using a temporal tolerance window of $\pm 10$ frames. Because attribute predictions are often noisy, this window accounts for the imprecision in localizing the exact switch points. We also report the total number of output tracklets to measure the overall cost of fragmentation.

\begin{itemize}
    \item \textbf{Precision}. This measures the accuracy of the splitter's cuts. It is the percentage of predicted splits that actually correspond to a real identity switch. A low precision score indicates that the splitter is creating too many false splits.

    \item \textbf{Recall}. This measures the coverage of the splitter. It is the percentage of ground truth identity switches that the splitter successfully caught. A low recall score means that the splitter is missing actual identity switches.

    \item \textbf{F1}. This is the harmonic mean of precision and recall. It provides a balanced summary score that evaluates the overall quality of the splitter without favoring over-splitting or under-splitting.
\end{itemize}


\subsubsection{Tracklet Connector}

HOTA \cite{hota} has become the primary evaluation metric across recent multi-object tracking literature. HOTA measures tracking quality by combining detection accuracy (DetA) and association accuracy (AssA) into a single score:

\begin{equation}
    HOTA = \sqrt{DetA \cdot AssA}
\end{equation}

Because we use the ground truth detections for all experiments, DetA is approximately constant for every method. Therefore the HOTA score mainly comes from AssA. AssA measures how correctly a tracker maintains the same identity for an object across frames. Our refinement framework is designed to improve its association accuracy, and is therefore reported separately to isolate the association improvements. All metrics are computed with the SoccerNet-TrackEval toolkit, a fork of the original TrackEval library \cite{trackeval}. 

\subsection{Baselines}\label{sec:baselines}
ALERT is compared against two baselines:

\begin{itemize}
    \item \textbf{Deep-EIoU without refinement}. The tracklets produced by the Deep-EIoU tracker on the  ground truth detections, without any post-processing.

    \item \textbf{Deep-EIoU + GTA}. The tracklets  produced by the Deep-EIoU tracker but using Global Tracklet Association as the offline tracklet refinement method. Comparing against GTA on identical Deep-EIoU tracklets isolates the contribution of the domain-specific attributes to refinement quality.
\end{itemize}

GTA is reproduced using the public GitHub implementation by Sun et al. \cite{gta}. 

\subsection{Training and Evaluation Setup}\label{sec:training_and_evaluation}





\section{Attribute Prediction}
Before evaluating the tracklet refinement components, we first examine the quality of the predicted attributes that serve as input signals. Since jersey and team splitters can only act on the information provided by their respective predictors, understanding where predictions fail is important to interpret splitter performance. We conduct a step-by-step analysis that tracks each ground truth  identity switch through four stages: whether the switch is actually present in the ground truth annotations, whether the predictor detects the relevant attribute on both sides of the switch, whether the predictions actually show two different values, and whether the splitter ultimately triggers a split. Table \ref{tab:attribute_analysis} summarizes the results across the 49 test sequences.

\

An important flaw of this  analysis is that the ground truth jersey numbers and team labels originate from the SoccerNet Game State Reconstruction dataset, which does not provide an attribute annotation per frame but per player identity. When we report that 724 switches involve players with ground truth jersey numbers, this means that the two players contain readable different jersey numbers somewhere in the match, not necessarily near the switch point. An identity switch that occurs for 5 frames will still have a ground truth jersey number even if that number can not be seen in those 5 frames, but somewhere later in the video. Consequently, the opportunity represents a theoretical ceiling based on player identity, not on visual opportunity. Similarly, the attributes are evaluated over the full tracklet segments before and after the switch point. Therefore, we do not use a frame window to analyze whether the splitter records a split near the identity switch, but whether it records a split at all.


\begin{table}[H]
\centering
\begin{tabular}{lrrrr}
\toprule
\textbf{Stage} & \multicolumn{2}{c}{\textbf{Jersey}} & \multicolumn{2}{c}{\textbf{Team}} \\
\cmidrule(lr){2-3} \cmidrule(lr){4-5}
 & Count & \% of prev. & Count & \% of prev. \\
\midrule
Total identity switches & 1048 & --- & 1048 & --- \\
Opportunity (GT attribute differs) & 724 & 69.1\% & 653 & 62.3\% \\
Detection (predictions on both sides) & 107 & 14.8\% & 643 & 98.5\% \\
Signal (predicted values differ) & 51 & 47.7\% & 293 & 45.6\% \\
Splitter fires & 13 & 25.5\% & 39 & 13.3\% \\
\bottomrule
\end{tabular}
\caption{Analysis of attribute prediction at identity switches.}
\label{tab:attribute_analysis}
\end{table}


\paragraph{Jersey Numbers}
Of the 1048 total ground truth identity switches, 724 (69.1\%) involve players wearing different jersey numbers. The OCR pipeline produces 107 confident predictions on both sides of the split point. This drop combines the cases where the jersey number was visible but the OCR pipeline failed to read it, and cases where no jersey was ever readable near the switch. The integrated legibility filter filters out almost all crops with no visible jersey numbers, however it sometimes also filters out crops where an jersey number is visible. 

\

Of these 107 switches where both sides have a jersey prediction, 51 produce a signal where the predicted numbers differ. The remaining cases represent instances where the OCR model predicts the same number for both players, either by confusion, occlusion, or players from the same team having the same jersey number. Of these 51 signals, the splitter records 13 (25.5\%) split points. The persistent filter requires a minimum number of  consistent  predictions within a lookahead window with a high agreement ratio and low entropy. This causes the splitter to not trigger when a jersey number is very briefly visible and therefore does not enable enough predictions. The remaining 38 signals are cases where the predictions do not meet the persistence threshold.

\

Further analysis shows that among the 13 splits that the jersey splitter triggers, the delay between the ground truth switch and the predicted split ranges from +5 to +236 frames (0.20s to 9.44s). As a result, the jersey splitter ensures that the tracking data remain pure over time, rather than catching identity switches the exact moment they occur.


\paragraph{Team Identities}
Team analysis reveals a fundamentally different bottleneck. Of 1048 identity switches, 653 involve a transition between players on different teams. Unlike jersey number prediction, the team classifier predicts a team identity on both sides of the split for 643 out of 653 switches. This is expected because team classification uses K-Means clustering to cluster players into 2 groups, which produces a prediction for every frame where a torso crops exists. The critical bottleneck for team predictions is signal accuracy. Of the 643 detected identity switches, only 293 produce a signal where the predicted teams on both sides differ.

\

The critical bottleneck for team prediction is prediction accuracy. Of the 643 detected switches, only 293 produce a signal where the predicted teams actually differ. Manual inspection reveals that the team classifier is highly reliable when players are clearly visible and not occluded, but highly sensitive to occlusion. Identity switches typically occur when two players are very close to each other, this causes the crop to contain the jersey of the occluding player rather than the player being tracked. In such cases, the classifiers correctly identifies the player visible but this does not match the player being tracked. This is not a classification error but a consequence of the input crops being contaminated by the occluding player. 

\

Of the 293 valid team signals, 39 (13.3\%) result in the splitter recording a split point. The persistence filter converts a smaller percentage of signals into splits compared to the jersey splitter, reflecting the noisier nature of team predictions.


\paragraph{Persistence Filter Sensitivity}
The persistence filter controls the conversion from signal to split shown in Table \ref{tab:attribute_analysis}. It controls how many consistent predictions are required, how far we look into the future predictions, and at what confidence or entropy threshold the predictions are considered valid. Tuning these values creates a tradeoff between missing ground truth identity switches and triggering false splits. To analyze this tradeoff, we use grid search to find the optimal parameter values. Figure \ref{fig:persistence_sensitivity}  shows the resulting precision versus recall for each configuration, where precision is defined as the fraction of all predicted splits that are true, and recall is the fraction of ground truth switches that are detected by at least one splitter. Note that these metrics differ from the intrinsic metrics in Table \ref{tab:splitter_results_final} in two ways: the tolerance window is the full video length instead of the 10 frames, and recall is computed only over switches where the relevant ground truth attributes differ (724 for jersey, 653 for team) rather than over all 1048 identity switches. The wider frame tolerance accounts for the fact that attribute splitters fire with delay (especially the jersey splitter), and the restricted recall computation excludes switches that are not detectable with the attribute splitters.

\

For the jersey splitter, precision ranges from 50\% to 92\% and recall from 0.1\% to 4.1\% across the parameter grid. The Pareto front reveals the precision-recall tradeoff: more aggressive settings increase recall at the cost of precision, primarily by predicting extra false splits. The false positives produced by the jersey splitter likely stem again from occlusion, which causes the crop to show a nearby player's jersey number, producing a persistent but incorrect signal.

For the team splitter, precision ranges from 52\% to 82\% and recall from 4.4\% to 12.6\%. The lookahead window is the most sensitive parameter, with shorter windows producing substantially more splits. The lower precision compared to jersey numbers reflects the noisier nature of team predictions in occluded regions, as discussed above.

\

The selected configurations are marked with stars in Figure \ref{fig:persistence_sensitivity} and lie on the Pareto front. These parameters have been evaluated on the validation set and the configuration with the highest HOTA score has been selected. These parameter values can be seen in Table \ref{tab:attribute_splitter_values}.

\begin{figure}[H]
\centering
\includegraphics[width=1\textwidth]{figures/persistence_precision_recall.png}
\caption{Persistence sensitivity}
\label{fig:persistence_sensitivity}
\end{figure}

\paragraph{Implications}
The attribute prediction analysis reveals two distinct failure modes. For jersey numbers, the system is bottlenecked by visibility and detection: the prediction pipeline predicts numbers on both sides of a switch only 14.8\% of the cases where the ground truth annotations indicate different jersey numbers. Improving detection coverage remains the most important work, whether through more aggressive detection strategies such as lowering the legibility thresholds or through richer input signals such as multiple camera angles. When the pipeline does produce a signal, 25.5\% of those signals lead do a split.

\

For team identity, the system is bottlenecked by occlusion sensitivity. The team classifier produces correct predictions when players are clearly visible, but identity switches occur during close player interactions where the occluding player contaminates the input crops. This is a fundamental limitation of appearance based team classification from a single camera view.
\

Both attribute analysis confirm that the persistence filters filter out most of the false signals while still converting a meaningful fraction of valid signals into actual splits. The persistence filter sensitivity analysis shows that the selected parameters lie on the Pareto front of the precision recall tradeoff, and that the system is robust to small parameter variations. 





\section{Tracklet Splitting}\label{sec:splitter_eval}
This section evaluates the tracklet splitting stage of the pipeline. Each splitting signal is assessed in isolation and in combination using the metrics described in Section \ref{sec:evaluation_methodology}, precision, recall, and F1. Table \ref{tab:splitter_results_final} presents the performance of each individual splitting signal evaluated against the baseline containing 1048 ground truth identity switches across 49 sequences. The results reveal a fundamental precision-recall trade-off, no single signal achieves both high precision and meaningful recall simultaneously.



\paragraph{Tolerance sensitivity}
The choice of the tolerance window has a substantial effect on the reported metrics in Table \ref{tab:splitter_results_final}. With a tolerance window of 150 frames (six seconds), the values change quite significantly: jersey precision doubles to 66.66\%, GTA changes from 290 true positives to 155 and the four signal configuration (STR + Jersey + Team + Bbox) registers 134 true positives instead of 108. A more permissive window risks counting predicted splits as correct when they were triggered by a different event  that happens to occur near a  ground truth  switch point within the same tracklet. For example when a ground truth switch at frame 200 occurs and the  splitter happens to split a tracklet at frame 340 for completely for wrong reasons, a tolerance of 150 frames would count that as a true positive even though the split had nothing to do with the identity switch at frame 200. A tolerance window of 10 frames is strict enough to ensure that a matched split is related to the identity switch it claims  to detect. 



\subsection{Individual Splitters}\label{subsec:individual_splitters}
This section addresses \ref{RSQ1} by evaluating which individual signal types are most effective at detecting identity switches.
\

\paragraph{Global Tracklet Association (GTA)}
Global Tracklet Association achieves the highest recall of any method at 14.79\% by predicting 2053 splits of which 155 are correct. However, as a trade-off it scores the lowest precision of all methods at 7.55\%. GTA aggressively over-splits: for every correct split, approximately thirteen false splits are predicted. This large amount of extra fragments puts a heavy demand on the downstream tracklet connector that has to reconnect these fragments. GTA serves as the current state-of-the-art baseline that we try to improve upon in this thesis.  


\paragraph{Spatio-Temporal ReID (STR)}
STR achieves the highest precision score of any splitting method at 80.65\%, predicting 93 splits of which 75 are correct. Despite its high precision, STR only scores 7.16\% for recall. This reflects the fact that the STR splitter only detects identity switches when a tracklet contains a temporal gap and when the ReID cosine distance between either side of the temporal gap exceeds a threshold. In broadcast football, the majority  of identity switches occur during crowded scenes such as corners or freekicks where players overlap each other continuously. These occurrences do not contain temporal gaps which causes STR to obtain a limited recall score. Despite the low recall, the STR splitter achieves the highest F1 score of all individual splitters with 13.15\%. The spatio-temporal ReID splitter is also computationally efficient: it only evaluates existing gaps in the tracking output and compares a small number of embeddings that are already computed by the tracker, making it negligible in cost relative to the attribute splitters.  

\paragraph{Jersey Number}
The jersey splitter predicts only 9 splits across all 49 test sequences, of these 9 predicted splits, all 9 are ground truth identity switches. However, 6 of the predicted splits are predicted outside of the 10 frames tolerance window, and are therefore counted as a false positive, causing the jersey splitter to only receive a 33.3\% precision score. The jersey splitter achieves the lowest recall score of all splitters at 0.29\%. The low recall reflects the fundamental visibility limitation of the jersey number splitter: jersey numbers are frequently unreadable due to occlusions, motion blur or low-resolution crops. Manual inspection of the 6 false positives revealed that all 6 correspond to genuine identity switches where the split point was delayed beyond the 10 frame tolerance. This occurs because jersey numbers are only detectable when the player's torso faces the camera, which may cause the splitter to delay its detection by many frames. The jersey splitter would achieve a 100\% precision rate, when relaxing the tolerance to 150+ frames. While the jersey number shows promise, the signal catches too few identity switches to be useful in isolation. 

\

Figure \ref{fig:pos_jersey_example} shows an example of how the jersey splitter correctly identified an identity switch by predicting two different persistent jersey numbers in the same tracklet. However, since the jersey numbers are not always visible, this splitter doesn't record the split point at the exact point that the switch occurred, instead it records the split point at the moment that the jersey number is first visible. In this example, because the second jersey number becomes visible just a  few frames after the identity switch occurred, the splitter performed effectively.

\begin{figure}[H]
\centering
\includegraphics[width=0.8\textwidth]{figures/jersey_splitter_pos_example.drawio.png}
\caption{Positive Jersey splitter example}
\label{fig:pos_jersey_example}
\end{figure}

Figure \ref{fig:neg_jersey_example} illustrates a tracklet containing an identity switch that remains unsplit for over 150 frames because no jersey number is visible. While the tracklet splitter functions as intended, the significant delay causes the evaluation metric to register a false positive. Despite this, the split still yields a positive effect on the HOTA score by successfully splitting an identity switch.

\begin{figure}[H]
\centering
\includegraphics[width=0.8\textwidth]{figures/jersey_splitter_bad_example.drawio.png}
\caption{Negative Jersey splitter example}
\label{fig:neg_jersey_example}
\end{figure}


\paragraph{Team Identity}
The team splitter significantly outperforms the jersey splitter, predicting 81 splits of which 37 are correct. The team splitter achieves a precision of 45.68\% and 3.53\%  recall. The team splitter achieves an F1 score of 6.55\% because of its relatively high average score for both precision and recall. Team predictions are more often available than jersey number predictions which causes the team splitter to predict almost ten times the amount of splits. However, team color is susceptible to lighting variation and partial occlusion, which explains the increased false positive rate.  Whereas the exact identity switch in Figure \ref{fig:neg_jersey_example} is hard for the jersey splitter to catch, the team splitter catches this identity switch instantly. In contrast, the team splitter makes more mistakes as seen in Figure \ref{fig:neg_team_example}, where because of occlusion the team splitter splits the tracklet, because of a wrong team prediction.

\paragraph{Bounding Box Anomaly}
The bounding box anomaly splitter detects abrupt velocity spikes of the bounding boxes and flags frames where the velocity threshold exceeds a certain threshold. The bounding-box splitter predicts 82 splits of which 39 are correct. The bounding box anomaly splitter achieves a precision of 47.56\% and a recall of 3.72\%. Its 39 correctly predicted splits include cases where neither the team, jersey, nor STR splitter triggers, making it a complementary addition. The lower precision than the team splitter indicates that in football legitimate rapid moving players can also produce velocity spikes. Despite similar performance, the bbox splitter is significantly more computationally efficient: it operates purely on the bounding-box trajectory already produced by the tracker and requires no additional attribute prediction, in contrast to the jersey and team splitters which require intensive inference pipelines. One limitation is that this evaluation uses ground-truth detections, where in practice noisy model detections may introduce additional velocity spikes that reduce precision.


\paragraph{Trajectory}
The trajectory splitter detects two players that overlap and suddenly change direction. The trajectory splitter predicts 143 splits of which only 7 are correct, achieving a precision of 4.9\% and a recall of 0.67\%. The high false-positive rate stems from the physical nature of football: tackles, dribbles and players crossing each other all produce the exact proximity and direction change pattern that the splitter is designed to detect, but without an actual identity switch. The trajectory splitter scores the second lowest F1 score of all splitters at 1.18\%. 


\begin{table}[ht]
    \centering

    \resizebox{\textwidth}{!}{%
    \begin{tabular}{lccccc}
        \toprule
        \textbf{Splitter Configuration} & \textbf{\# Tracklets} & \textbf{Pred. Spl.} & \textbf{Recall $\uparrow$} & \textbf{Precision $\uparrow$} & \textbf{F1 $\uparrow$} \\
        \midrule
        Baseline                            & 1725 & 0    & 0.0000 & ---    & ---    \\
        GTA                                 & 3776 & 2053 & \textbf{0.1479} & 0.0755 & 0.1 \\
        Spatio Temporal Re-ID (STR)         & 1818 & 93   & 0.0716 & \textbf{0.8065} & \textbf{0.1315} \\
        Jersey                              & 1734 & 9    & 0.0029 & 0.3333 & 0.0057 \\
        Team                                & 1805 & 81   & 0.0353 & 0.4568 & 0.0655 \\
        Bbox                                & 1807 & 82   & 0.0372 & 0.4756 & 0.0690 \\
        Trajectory                          & 1868 & 143  & 0.0067 & 0.0490 & 0.0118 \\
        \midrule
        Jersey + Team                       & 1814 & 90   & 0.0382 & 0.4444 & 0.0703 \\
        Jersey + Team + Bbox                & 1880 & 156  & 0.0620 & 0.4167 & 0.1080 \\
        Jersey + Team + Bbox + Trajectory   & 2024 & 300  & 0.0687 & 0.2400 & 0.1068 \\
        STR + Jersey + Team + Bbox          & 1929 & 206  & 0.1031 & \textbf{0.5243} & \textbf{0.1722} \\
        STR + Jersey + Team + Bbox + Trajectory & 2071 & 348 & \textbf{0.1097} & 0.3305 & 0.1648 \\
        \bottomrule
    \end{tabular}}
    \caption{Performance comparison of tracklet splitter configurations. All methods are evaluated against a ground truth of 1,048 identity switches.}
    \label{tab:splitter_results_final}
\end{table}

\subsection{Splitter Combinations and Ordering}
The previous section evaluated each splitting signal in isolation. This section addresses the second part of \ref{RSQ1} by investigating how combining multiple splitters affects the splitting performance. The lower half of Table \ref{tab:splitter_results_final} summarizes every configuration that was evaluated. 

\paragraph{Attribute combinations}
Combining jersey number and team identities increases the F1 score to 7.03\% at 44.44\% precision and a recall of 3.82\%. The jersey splitter predicts 9 additional splits not detected by the team splitter, adding 3 additional true positives and 6 false positives on top of the team splitter alone. While the jersey splitter alone is too sparse to be useful in isolation, it provides a small but consistent improvement when used as a complementary splitter.

\

Adding the bounding box anomaly splitter on top of the jersey and team splitter configuration raises the F1 score to 10.80\%, with recall increasing to 6.2\%. Precision drops to 41.67\% due to the noisier geometric signal. The bounding  box splitter contributes 25 additional true positives, indicating that it detects a complementary class of switches: fast spatial displacements that occur without a change in visible jersey number or team identity. 


\paragraph{Integrating Spatio-Temporal ReID}
Adding STR on top of the Jersey + Team + Bbox configuration produces the largest F1 gain of any individual addition, improving F1 score from 10.80\% to 17.22\%. STR adds 50 additional predicted splits of which 43 are ground truth identity switches not detected by the other splitters. This high marginal precision of 86.02\% is substantially above the other splitters, raising the precision from 41.67\% to  52.43\%, making STR the most efficient single addition.

\paragraph{Trajectory as fifth splitter}
Adding the trajectory splitter on top of the best performing  STR + Jersey + Team + Bbox combination achieves the second highest F1 score of 16.48\% by substantially increasing recall, but reducing precision from 52.43\% to 33.05\%. Since trajectory alone obtains a very low F1 score by having a very low precision, it only improves the F1 score by adding recall. However, because the trajectory splitter introduces many false positives, the tracklet connector has a harder job connecting all the tracklets back together.  

\paragraph{Takeaway}
Each splitter captures a distinct subset of identity switches, confirming that attribute, geometric, and ReID signals are complementary. The domain specific attributes (jersey numbers and team identities) and STR contribute the highest precision, while the geometric signals (bounding box and trajectory) contribute to a higher recall at lower precision. The recommended  four signal configuration (STR + Jersey + Team + Bbox) achieves the highest F1 score of 17.22\% at 52.43\% precision capturing 108 of the 1048 identity switches while predicting 98 false splits.

% training graphs


\subsection{Parameter Sensitivity}
The previous sections evaluated the splitter's performance and configurations using fixed parameter values. These values were selected using grid search on the validation set. This section evaluates the robustness and sensitivity of the parameters of each splitter. For every parameter, we test a grid of values while keeping other parameters constant. Each configuration is evaluated through the full pipeline and the resulting HOTA score on the validation set is recorded. Figure \ref{fig:sensitivity_overview} provides an overview of the HOTA scores spread across all configurations for each splitter. The detailed sensitivity plots are shown in Appendix \ref{app:sensitivity_plots}.

\begin{figure}[H]
\centering
\includegraphics[width=1\textwidth]{figures/splitter_sensitivity_overview.png}
\caption{Negative Jersey splitter example}
\label{fig:sensitivity_overview}
\end{figure}


\paragraph{Jersey Splitter}
The jersey splitter is tuned over four parameters: entropy threshold, minimum persistence count, lookahead window, and persistence ratio. The total HOTA range is 0.305 points from 88.680 to 88.985. The most sensitive parameter is the minimum persistence count, which produces a HOTA range of 0.299. Using a low persistence count, such as 10 predictions, causes the splitter to trigger on brief, unreliable jersey predictions. The entropy threshold, lookahead window, and persistence ratio each produce ranges below 0.06, indicating that the jersey splitter parameters are robust to small variations around the selected values.

\paragraph{Team Splitter}

\paragraph{Spatio-Temporal ReID Splitter}


\paragraph{Bounding Box Anomaly Splitter}
The bounding box anomaly splitter relies on five parameters that control how velocity spikes are detected. To assess which parameters have the most influence on performance, and how robust the parameters are against changes, we conduct a sensitivity analysis. The results are evaluated  using HOTA on the validation set with only the bounding box splitter active.

\


\paragraph{Trajectory Splitter}
The trajectory splitter is tuned over five parameters: proximity distance, velocity window, minimum velocity change, direction change threshold, and swap similarity threshold. Across 25 configurations, the total HOTA range is 1.658 points from 86.445 to 88.103, the largest of all splitters. The direction change threshold is the most sensitive parameter with a range of 1.643. HOTA increases significantly when increasing the threshold from 60 to 180 degrees. The swap similarity is the second most sensitive parameter with a range of 0.94, with HOTA decreasing when the threshold is lowered from 0.95 to 0.7. The velocity window, proximity distance, and minimum velocity show moderate sensitivity. 

\

Crucially, the configurations that achieve the highest HOTA scores are those where the trajectory splitter produces zero splits in all sequences. In these configurations, the trajectory splitter is effectively disabled, and the HOTA score equals the XGBoost connected tracklets without splitting applied beforehand. This means that all configurations where the trajectory splitter actually triggers produces a lower HOTA score than not splitting at all. The sensitivity analysis confirms what the intrinsic evaluation in Section \ref{subsec:individual_splitters} suggested: the geometric signals used by the trajectory splitter are too noisy for broadcast football footage. Therefore, the trajectory splitter is excluded from the recommended pipeline configuration. 



\subsection{Analysis}
\paragraph{Precision-Recall Tradeoff}
Table \ref{tab:splitter_results_final} reveals a consistent pattern across all splitter configurations: a higher recall comes with a lower precision. This tradeoff is important to understand, because it has a direct consequence for the final tracking quality. Every false split creates an additional fragment that the tracklet connector has to process and when the tracklet splitter creates too many fragments, it is more likely that the tracklet connector fails to reconnect fragments that belong together, ultimately reducing the final tracking quality. 

\

This effect  is clearly visible when comparing the four signal configuration (STR + Jersey + Team + Bbox) against the five signal configuration that adds the trajectory splitter. Although the five signal configuration achieves a high F1 score of 16.48\%, it achieves one of the lowest precision at 33.05\%. The trajectory splitter adds 142 extra fragments, of which 135 are false positives. When both configurations are evaluated through the entire pipeline using a fixed XGBoost merger, the five signal configuration achieves a 0.8 points HOTA score lower than the four signal  configuration. The trajectory splitter's additional true positives do not compensate for the large amount of false positives.

\

This result demonstrates that optimizing the splitter for F1 score does not necessarily translate to the best end-to-end tracking performance. For this reason, Jersey + Team + Bbox + STR is selected as the recommended splitter configuration for all subsequent experiments, as it provides the best balance between recall and precision for the tracklet connector.


\paragraph{The Recall Ceiling}
Across all the splitter combinations, the most aggressive configuration recovers only 14.79\% of the 1048 ground truth identity switches. The remaining 85.21\% of identity switches go undetected because of football fundamentals and the design of the tracklet splitter.

\

The majority of uncaught identity switches occur during  continuous tracklets where two players on the same team are involved. In these cases the STR splitter is ineffective because no temporal gap exists, the team identity does not change because both players wear the same uniform, and the jersey number is often not visible because of occlusion or player torsos facing away from the camera. The bounding box anomaly splitter can sometimes detect the velocity spikes caused by an identity switch, but when two players move very close to each other the velocity spike is often too small to exceed the threshold.  

\ 

A second class of uncaught switches involves identity switches that occur over a very short window. The attribute splitters require a persistent attribute change before triggering a split. This causes brief identity switches to resolve before the splitters can detect them. Lowering these persistence thresholds could catch these switches, however would substantially increase the amount of false splits.

\

A third class consists of switches that occur during rapid camera movement. This causes ReID embeddings and attribute  predictions to be of very low quality and often remain unusable. In these frames, no signal can be utilized for detecting the identity switch.

\

Ultimately, the low recall reflects a fundamental limitation of the chosen splitting signals rather than a failure of the splitting framework itself. The attribute and geometric splitters are most effective when an identity switch produces a visible, persistent change in at least one of the attributes. Identity switches between players on the same team remain the hardest case and would likely require richer signals such as pose estimation, action recognition, or multiple camera views to resolve consistently. 


\subsection{Qualitative Analysis}

% hota table

% SHAP importance plot
% ablation study tracklet splits

\begin{figure}[h!]
\centering
\includegraphics[width=1\textwidth]{figures/duration_distribution.png}
\caption{Identity switch duration analysis}
\label{fig:duration_analysis}
\end{figure}

\begin{figure}[h!]
\centering
\includegraphics[width=1\textwidth]{figures/gap_and_distance.png}
\caption{Identity switch duration analysis}
\label{fig:duration_analysis}
\end{figure}

\begin{figure}[h!]
\centering
\includegraphics[width=1\textwidth]{figures/summary_table.png}
\caption{Identity switch duration analysis}
\label{fig:duration_analysis}
\end{figure}


\section{Tracklet Connecting}\label{sec:tracklet_connecting}

This section evaluates the tracklet connecting stage of the ALERT framework. After the splitter produces purified tracklet fragments, the connector must decide which fragments belong to the same player and merge them back into complete trajectories. We evaluate three connector architectures of increasing complexity: a rule-based connector that uses hand-crafted heuristics (Section~\ref{subsec:rule_based_connector}), an XGBoost-based connector that learns merge decisions from pairwise features (Section~\ref{subsec:xgboost_connector}), and a transformer-based connector that jointly reasons over all fragments (Section~\ref{subsec:transformer_connector}). All connectors receive the same splitter output (STR + Jersey + Team + Bbox) and are evaluated using HOTA on the same dataset splits. The XGBoost connector is the most extensively studied, with experiments covering training data generation, merge threshold tuning, hard constraint ablation, and feature ablation.


\subsection{Rule-based Connector}\label{subsec:rule_based_connector}


\subsection{XGBoost Connector}\label{subsec:xgboost_connector}

The XGBoost connector frames tracklet merging as a binary classification problem. For each pair of tracklet fragments, the model receives 46 pairwise features derived from six attribute groups (jersey numbers, team identity, temporal proximity, spatial proximity, ReID embeddings, and SigLIP embeddings) and predicts the probability that the two fragments belong to the same player. These probabilities are converted into distances ($d = 1 - p$) and fed into hierarchical agglomerative clustering with average linkage. Fragment pairs whose distance falls below a merge threshold are merged. Before classification, three hard constraints can override the model prediction: temporally overlapping fragments, fragments with conflicting confident jersey numbers, and fragments with conflicting consistent team labels are assigned a distance of 2.0, preventing them from ever being merged regardless of the model output.

\

The XGBoost connector involves several design choices that affect end-to-end tracking quality: how training data is generated, what merge distance threshold is used, whether hard constraints help or hurt, and which features contribute most to performance. The following paragraphs investigate each of these choices in turn.


\paragraph{Data Generation and Training}\label{par:data_gen}

The XGBoost classifier requires labelled training pairs of tracklet fragments, where each pair is labelled as positive (same player) or negative (different players). Generating this training data involves several choices that directly influence the distribution of training examples and, consequently, the quality of the learned merge function. We identify four key parameters that control the training data generation process:

\begin{itemize}
    \item \textbf{Negative sampling ratio}. Because the number of possible negative pairs grows quadratically with the number of fragments while positive pairs are comparatively rare, using all negative pairs produces a heavily imbalanced dataset. The negative sampling ratio controls how many negative pairs are sampled per positive pair. We test ratios of 3, 5, and 10, as well as no subsampling (using all negative pairs).

    \item \textbf{Tracklet purity threshold}. Each training fragment is assigned a purity score based on the fraction of detections that belong to the majority ground truth identity. Fragments below the purity threshold are excluded from training. Lower thresholds include noisier fragments that better reflect the impure fragments produced by the splitter at inference time. We test thresholds of 0.6, 0.8, and 1.0.

    \item \textbf{Minimum tracklet length}. Short tracklet fragments carry less information for computing reliable features, particularly for attributes that require temporal aggregation such as team consistency or jersey entropy. The minimum length filter excludes fragments shorter than a given number of frames from the training set. We test minimum lengths of 0, 5, and 10 frames.

    \item \textbf{Splitter application}. Training data can be generated from the original tracker output or from the output of the tracklet splitter. When the splitter is applied, the training fragments more closely resemble the fragments the model will encounter at inference time. We test both settings.
\end{itemize}

The full grid of $4 \times 3 \times 3 \times 2 = 72$ configurations is trained using identical XGBoost hyperparameters, and each model is evaluated through the complete pipeline (merge, write MOT files, compute HOTA) on both the validation and test splits.

\paragraph{Validation results}
Table~\ref{tab:xgboost_top10_val} presents the ten configurations that achieve the highest HOTA scores on the validation set. All top-10 configurations use a negative sampling ratio of 3 or 5 and a purity threshold of 0.6 or 0.8. The best configuration, \texttt{neg3\_minlen0\_purity0.6\_splitTrue}, achieves a validation HOTA of 90.257 with an AssA of 86.628. The spread across the top 10 is only 0.150 HOTA points, indicating that the best configurations perform similarly.

\begin{table}[H]
\centering
\resizebox{\textwidth}{!}{%
\begin{tabular}{lcccccc}
\toprule
\textbf{Configuration} & \textbf{Val HOTA $\uparrow$} & \textbf{Val AssA} & \textbf{Test HOTA} & \textbf{Test AssA} & \textbf{Val F1} \\
\midrule
neg3\_minlen0\_purity0.6\_splitTrue   & \textbf{90.257} & 86.628 & 89.666 & 85.926 & 0.770 \\
neg5\_minlen5\_purity0.8\_splitTrue   & 90.202 & 86.502 & 89.799 & 86.183 & 0.743 \\
neg5\_minlen5\_purity0.6\_splitFalse  & 90.197 & 86.495 & 89.729 & 86.031 & 0.751 \\
neg5\_minlen0\_purity0.8\_splitFalse  & 90.192 & 86.496 & 89.559 & 85.710 & 0.735 \\
neg5\_minlen0\_purity0.6\_splitFalse  & 90.170 & 86.454 & 89.603 & 85.800 & 0.722 \\
neg3\_minlen0\_purity0.6\_splitFalse  & 90.146 & 86.423 & 89.665 & 85.914 & 0.763 \\
neg3\_minlen5\_purity0.6\_splitTrue   & 90.141 & 86.406 & 89.918 & 86.403 & 0.784 \\
neg5\_minlen0\_purity0.8\_splitTrue   & 90.132 & 86.367 & 89.483 & 85.578 & 0.720 \\
neg3\_minlen0\_purity0.8\_splitTrue   & 90.123 & 86.369 & 89.936 & 86.444 & 0.769 \\
neg5\_minlen10\_purity0.6\_splitFalse & 90.107 & 86.322 & 89.758 & 86.078 & 0.755 \\
\bottomrule
\end{tabular}}
\caption{Top-10 XGBoost connector configurations ranked by validation HOTA. All models use a default merge threshold of 0.5 and average linkage. Val F1 is the binary classification F1 on held-out validation pairs.}
\label{tab:xgboost_top10_val}
\end{table}

\paragraph{Parameter analysis}
To understand which data generation parameters matter most, Table~\ref{tab:xgboost_param_trends} reports the mean validation HOTA for each parameter level, averaged over all other parameters.

\begin{table}[H]
\centering
\begin{tabular}{llccc}
\toprule
\textbf{Parameter} & \textbf{Value} & \textbf{Mean HOTA $\uparrow$} & \textbf{Std} & \textbf{Range} \\
\midrule
\multirow{4}{*}{Negative ratio}
 & None       & 88.926 & 0.964 & [87.473, 89.741] \\
 & 3          & 89.757 & 0.421 & [88.998, 90.257] \\
 & 5          & \textbf{89.806} & 0.422 & [89.026, 90.202] \\
 & 10         & 89.465 & 0.450 & [88.669, 90.000] \\
\midrule
\multirow{3}{*}{Purity threshold}
 & 0.6        & \textbf{89.918} & 0.201 & [89.607, 90.257] \\
 & 0.8        & 89.760 & 0.328 & [88.839, 90.202] \\
 & 1.0        & 88.786 & 0.743 & [87.473, 89.851] \\
\midrule
\multirow{3}{*}{Min. tracklet length}
 & 0          & \textbf{89.536} & 0.741 & [87.473, 90.257] \\
 & 5          & 89.484 & 0.698 & [87.635, 90.202] \\
 & 10         & 89.445 & 0.666 & [87.682, 90.107] \\
\midrule
\multirow{2}{*}{Splitter applied}
 & True       & \textbf{89.576} & 0.678 & [87.595, 90.257] \\
 & False      & 89.400 & 0.707 & [87.473, 90.197] \\
\bottomrule
\end{tabular}
\caption{Mean validation HOTA per parameter level, averaged over all other parameters. Bold indicates the best level per parameter.}
\label{tab:xgboost_param_trends}
\end{table}

The negative sampling ratio is the most influential parameter. Using all negative pairs (ratio = None) produces the lowest mean HOTA of 88.926 and the highest variance ($\sigma = 0.964$), because the resulting class imbalance causes the model to underpredict merges. Ratios of 3 and 5 perform best (89.757 and 89.806 respectively), providing enough negative examples for the model to learn discriminative features without overwhelming the positive class. A ratio of 10 slightly degrades performance, suggesting that the optimal balance lies in the range of 3 to 5 negatives per positive.

\

The purity threshold has the second largest effect. Configurations with purity $= 1.0$ score 1.13 HOTA points below those with purity $= 0.6$ on average. Training exclusively on pure fragments creates a distribution mismatch: at inference time, the splitter produces fragments that are not perfectly pure, and the model has never seen feature vectors computed from impure fragments. Lowering the purity threshold to 0.6 exposes the model to noisier training examples that better reflect the inference distribution, acting as a form of implicit regularization. This finding is consistent across all negative ratios and minimum lengths.

\

The minimum tracklet length has a negligible effect on mean HOTA, with differences below 0.1 points across all three levels. This suggests that the feature computation is robust to short fragments, and that excluding them from training does not meaningfully change the learned decision boundary.

\

Applying the splitter before generating training data provides a small but consistent benefit of 0.176 HOTA points on average. This confirms that training on splitter output produces features that are more representative of the inference-time distribution, where the model operates on split fragments rather than raw tracker output.


\paragraph{Classifier metrics versus end-to-end HOTA}
An important finding of this experiment is that standard binary classification metrics are poor proxies for end-to-end tracking quality. The Spearman rank correlation between validation F1 and validation HOTA across the 72 configurations is $\rho = 0.632$, and between validation AUC and validation HOTA it is only $\rho = 0.497$. This means that selecting a model based on the highest classifier F1 or AUC does not reliably identify the model that produces the best tracking output. The disconnect arises because the classifier operates on individual pairs, while HOTA evaluates the global assignment after clustering. A model with a slightly lower F1 may produce a distance matrix that is better suited for hierarchical clustering, for instance by producing well-separated clusters even if its per-pair accuracy is marginally lower.

\

This finding has a practical consequence for model selection: the XGBoost connector must be evaluated through the full pipeline (merge, write MOT files, compute HOTA) rather than relying on classifier metrics as a proxy. All subsequent experiments therefore use validation HOTA as the selection criterion.


\paragraph{Validation versus test set agreement}
To verify that validation performance generalises, we evaluate all 72 configurations on both the validation and test splits. The Spearman rank correlation between validation and test HOTA is $\rho = 0.870$ ($p < 10^{-22}$) and the Pearson correlation is $r = 0.949$, indicating strong agreement. However, only 4 of the top-10 validation configurations also appear in the top-10 test configurations. The best test configuration (\texttt{neg3\_minlen10\_purity0.8\_splitTrue}, HOTA = 89.978) ranks outside the validation top 10. This partial overlap is expected: with 72 tightly clustered configurations, small per-sequence fluctuations can reorder the ranking even when the overall correlation is high.

\

The parameter-level trends are consistent across both splits: negative ratios of 3 and 5 outperform 10 and None, purity $= 1.0$ is consistently the worst, and the splitter application provides a small benefit. This consistency confirms that the validation set provides a reliable signal for parameter selection, even if the exact ranking of individual configurations shifts.


\paragraph{Selected configuration}
Based on the validation HOTA ranking, we select \texttt{neg3\_minlen0\_purity0.6\_splitTrue} as the base XGBoost model for all subsequent experiments (merge threshold sweep, hard constraint ablation, feature ablation). This configuration achieves a validation HOTA of 90.257 and a test HOTA of 89.666. Although it is not the top-performing configuration on the test set, selecting by validation performance is the methodologically correct approach as it prevents implicit data leakage from the test set into the model selection process.


\paragraph{Hard Constraints}\label{par:hard_constraints}


\paragraph{Feature Ablation}\label{par:feature_ablation}


\subsection{Transformer Connector}\label{subsec:transformer_connector}
